%%%%%%%%%%%%%%%%%%%%%%%%%%%%%%%%%%%%%%%%%%%%%%%%%%%%%%%%%%%%%%%%%%%%%%%%%%%%%%%%
% AMS Beamer series / Silver Blue / Template
% Andrea Omicini
% Alma Mater Studiorum - Università di Bologna
% mailto:andrea.omicini@unibo.it
%
% This template is released into the public domain under CC0 1.0 Universal:
% use it, adapt it and build on it freely, with no obligation to credit.
% See the LICENSE file, or https://creativecommons.org/publicdomain/zero/1.0/
%
% The AMSSilverBlue style it uses is covered by its own licence (LPPL 1.3c):
% https://github.com/andreaomicini/beamer-AMSSilverBlue
%%%%%%%%%%%%%%%%%%%%%%%%%%%%%%%%%%%%%%%%%%%%%%%%%%%%%%%%%%%%%%%%%%%%%%%%%%%%%%%%
\documentclass[presentation,cesena,apice]{beamer}\mode<presentation>{\usetheme{AMSSilverBlue}}
%\documentclass[presentation]{beamer}\mode<presentation>{\usetheme{AMSSilverBlue}}
%\documentclass[handout]{beamer}\mode<handout>{\usetheme{AMSSilverBlue}}
%%%%%%%%%%%%%%%%%%%%%%%%%%%%%%%%%%%%%%%%%%%%%%%%%%%%%%%%%%%%%%%%%%%%%%%%%%%%%%%%
\def\BibTeX{{\rm B\kern-.05em{\sc i\kern-.025em b}\kern-.08em
    T\kern-.1667em\lower.7ex\hbox{E}\kern-.125emX}}
\renewcommand{\thefootnote}{\fnsymbol{footnote}}
% a colour swatch, framed so the pale ones still read against white
\newcommand{\swatch}[1]{\fcolorbox{sbgreyline}{#1}{\rule{0pt}{1.3ex}\rule{2.2ex}{0pt}}}
% a group heading inside the palette tables
\newcommand{\swgroup}[1]{\multicolumn{2}{@{}l}{\textcolor{sbbluedeep}{\textbf{#1}}}}
% version -- Major.Minor[.Patch]; the release time-stamp is appended by the CI
\newcommand{\templatemajor}{1}
\newcommand{\templateminor}{0}
% \templatepatch is optional: comment the next line out to release as Major.Minor
\newcommand{\templatepatch}{1}
\ifdefined\templatepatch{}
  \newcommand{\templateversion}{\templatemajor.\templateminor.\templatepatch}
\else
  \newcommand{\templateversion}{\templatemajor.\templateminor}
\fi
%%%%%%%%%%%%%%%%%%%%%%%%%%%%%%%%%%%%%%%%%%%%%%%%%%%%%%%%%%%%%%%%%%%%%%%%%%%%%%%%
\title[AMS Silver Blue Beamer Style]
{AMS Silver Blue Beamer Style}
%
\subtitle[AMS Series Templates]
{AMS Series Templates}
%
\author[\sspeaker{Omicini} \and Else]
{\speaker{Andrea Omicini} \and Someone Else}
%
\institute[DISI, Univ. Bologna]
{\url{mailto:andrea.omicini@unibo.it} \qquad \url{mailto:someone.else@unibo.it}\\
Dipartimento di Informatica \textendash{} Scienza e Ingegneria (DISI)\\
\textsc{Alma Mater Studiorum} \textendash{} Università di Bologna}
%
\date[v.\ \templateversion]{version \templateversion}
%
%%%%%%%%%%%%%%%%%%%%%%%%%%%%%%%%%%%%%%%%%%%%%%%%%%%%%%%%%%%%%%%%%%%%%%%%%%%%%%%%
\begin{document}
%%%%%%%%%%%%%%%%%%%%%%%%%%%%%%%%%%%%%%%%%%%%%%%%%%%%%%%%%%%%%%%%%%%%%%%%%%%%%%%%

%/////////
\frame{\titlepage}
%/////////

%%===============================================================================
%\section*{Outline}
%%===============================================================================
%
%%/////////
%\frame[c]{\tableofcontents[hideallsubsections]}
%%/////////

%===============================================================================
\section{Premises}
%===============================================================================

%/////////
\begin{frame}[c]{Academic Presentations}
%
\begin{itemize}
	\item technical presentations are common in the academic environment
	\item \alert{\LaTeX{}}\footnote{\uurl{https://www.latex-project.org}} is typically used for scientific literature
	\begin{itemize}
		\item for better formulas, bibliography, typography, \ldots
	\end{itemize}
	\item so, \LaTeX{} is the most typical \emph{toolbox} for scientific presentations
	\begin{itemize}
		\item which easily translates into teaching presentations, too
	\end{itemize}
\end{itemize}
%
\end{frame}
%/////////

%/////////
\begin{frame}[c]{Beamer}
%
\begin{itemize}
	\item \alert{Beamer} is a \LaTeX{} document class for creating presentation \emph{slides}, with a wide range of \emph{templates} and a set of features for making \emph{slideshow effects}
	\begin{description}
		\item[repo] \uurl{https://github.com/josephwright/beamer}
		\item[CTAN] \uurl{https://ctan.org/pkg/beamer}
	\end{description}
	\item even though other styles for \LaTeX{} presentations exist, Beamer is by far the best way today
\end{itemize}
%
\end{frame}
%/////////

%/////////
\begin{frame}[c]{AMS}
%
\begin{description}
	\item[AMS] Alma Mater Studiorum\textendash{}Università di Bologna
	\item[Silver Blue] the colours of the old \alert{AMSCesenaBleu} style, recovered
	\begin{itemize}
		\item its \textcolor{sbblue}{\textit{French Blue}}, over the \textcolor{sbgreymid}{\textit{silver}} greys it used to take from Beamer's own \texttt{beaver}
		\item since the \emph{Campus di Cesena} is where the style originated, its seal can be used for watermarking every slide
	\end{itemize}
\end{description}
%
\end{frame}
%/////////

%/////////
\begin{frame}[c]{Aim}
%
\begin{itemize}
	\item many occasions for presentation
	\item many different people in the group
	\item a common reference could be useful
	\item hence, this style
	\item in the colours of the old \alert{AMSCesenaBleu} style, recovered
	\begin{itemize}
		\item the same style as \alert{AMSUniBo} in everything else---options, commands, bibliography
		\item but the faces, which stay \LaTeX{}'s own Computer Modern
	\end{itemize}
\end{itemize}
%
\end{frame}
%/////////

%/////////
\begin{frame}[c]{Who}
%
\begin{itemize}
	\item mostly, myself and my collaborators
	\item everybody else, be my guest
\end{itemize}
%
\end{frame}
%/////////

%/////////
\begin{frame}[c]{Distribution}
%
\begin{block}{Public on GitHub}
	\begin{itemize}
		\item the style is distributed on \uurl{https://github.com/andreaomicini/beamer-AMSSilverBlue/}
		\item the releases are available at \uurl{https://github.com/andreaomicini/beamer-AMSSilverBlue/releases}
		\item this template is distributed on \uurl{https://github.com/andreaomicini/beamer-AMSSilverBlue-template/}
	\end{itemize}
\end{block}
%
\end{frame}
%/////////

%===============================================================================
\section{Use}
%===============================================================================

%/////////
\begin{frame}[c,fragile]{Declaration}
%
\begin{block}{\texttt{\textbackslash{}documentclass}}
Your main Beamer \texttt{.tex} file should open with the declaration
%
\begin{verbatim}
    \documentclass[presentation]{beamer}
        \mode<presentation>{\usetheme{AMSSilverBlue}}
\end{verbatim}
%
so as to use the AMSSilverBlue Beamer style
\end{block}
%
\end{frame}
%/////////

%/////////
\begin{frame}[c,fragile]{Bibliography Style}
%
\begin{block}{\texttt{apalike-AMS}}
Your main Beamer \texttt{.tex} file should include the declaration
\begin{verbatim}
    \bibliographystyle{apalike-AMS}
\end{verbatim}
so as to use the AMSSilverBlue \BibTeX{} style
\end{block}
%
\end{frame}
%/////////

%/////////
\begin{frame}[c,fragile]{Template}
%
\begin{block}{\texttt{AMSSilverBlue-template.tex}}
This template's sources can be used as a simple example of how to use this Beamer style
\end{block}
%
\end{frame}
%/////////


%===============================================================================
\section{Style}
%===============================================================================

%/////////
\begin{frame}[c,fragile]{Colours}
%
The palette is the old \alert{AMSCesenaBleu}'s, and every colour takes the role its
counterpart has in \alert{AMSUniBo}: \textcolor{sbblue}{\textit{blue}} identifies,
\textcolor{sbbluedeep}{\textit{deep blue}} and \textcolor{sbgreymid}{\textit{silver}} structure
%
\begin{exampleblock}{The four that carry the style}
\begin{tabular}{@{}llll@{}}
\swatch{sbblue}     & \texttt{sbblue}     & \texttt{\#0073BA} & head line, author, bullets, alerts, links \\
\swatch{sbbluedeep} & \texttt{sbbluedeep} & \texttt{\#003F66} & block headers, emphasis, \textsc{toc} \\
\swatch{sbsilver}   & \texttt{sbsilver}   & \texttt{\#B3B3B3} & separation lines, and the bands and fills \\
\swatch{sbgreymid}  & \texttt{sbgreymid}  & \texttt{\#5A5A5A} & subtitle, date, footnotes \\
\end{tabular}
\end{exampleblock}
%
\begin{itemize}
	\item \texttt{sbblue} is \emph{French Blue}, the one colour \alert{AMSCesenaBleu} defined;
	      \texttt{sbbluedeep} is the shade it gave its block headers
	\item \texttt{sbsilver} and its tints are the greys of Beamer's \texttt{beaver}, which
	      \alert{AMSCesenaBleu} was built on
\end{itemize}
%
\end{frame}
%/////////

%/////////
\begin{frame}[c]{The Full Palette}
%
\footnotesize
\begin{columns}[T]
\begin{column}{.33\textwidth}
\begin{tabular}{@{}ll@{}}
\swgroup{identity}\\
\swatch{sbblue} & \texttt{sbblue} \\
\swatch{sbbluedark} & \texttt{sbbluedark} \\
\end{tabular}
\medskip
\begin{tabular}{@{}ll@{}}
\swgroup{structure}\\
\swatch{sbbluedeep} & \texttt{sbbluedeep} \\
\swatch{sbbluelight} & \texttt{sbbluelight} \\
\swatch{sbbluetint} & \texttt{sbbluetint} \\
\swatch{sbsilver} & \texttt{sbsilver} \\
\swatch{sbsilvertint} & \texttt{sbsilvertint} \\
\swatch{sbsilvertint2} & \texttt{sbsilvertint2} \\
\end{tabular}
\end{column}
\begin{column}{.33\textwidth}
\begin{tabular}{@{}ll@{}}
\swgroup{text and neutrals}\\
\swatch{sbblack} & \texttt{sbblack} \\
\swatch{sbtext} & \texttt{sbtext} \\
\swatch{sbgrey} & \texttt{sbgrey} \\
\swatch{sbgreymid} & \texttt{sbgreymid} \\
\swatch{sbgreyline} & \texttt{sbgreyline} \\
\swatch{sbgreybg} & \texttt{sbgreybg} \\
\swatch{sbwhite} & \texttt{sbwhite} \\
\end{tabular}
\end{column}
\begin{column}{.33\textwidth}
\begin{tabular}{@{}ll@{}}
\swgroup{block fills}\\
\swatch{sbbluefill} & \texttt{sbbluefill} \\
\swatch{sbblueborder} & \texttt{sbblueborder} \\
\swatch{sbsilverfill} & \texttt{sbsilverfill} \\
\swatch{sbsilverborder} & \texttt{sbsilverborder} \\
\end{tabular}
\end{column}
\end{columns}
%
\bigskip
Each colour is a shade \alert{AMSCesenaBleu} or \texttt{beaver} actually used, where one
existed; the rest are guessed from them, and say so in \texttt{beamercolorthemesilverblue.sty}
%
\end{frame}
%/////////

%/////////
\begin{frame}[c,fragile]{Blocks}
%
\begin{block}{This is a \texttt{block} environment}
\begin{verbatim}
\begin{block}{...}
...
\end{block}
\end{verbatim}
\end{block}
%
\begin{exampleblock}{This is an \texttt{exampleblock} environment}
\begin{verbatim}
\begin{exampleblock}{...}
...
\end{exampleblock}
\end{verbatim}
\end{exampleblock}
%
\begin{alertblock}{This is an \texttt{alertblock} environment}
\begin{verbatim}
\begin{alertblock}{...}
...
\end{alertblock}
\end{verbatim}
\end{alertblock}
%
\end{frame}
%/////////

%===============================================================================
\section{New Commands}
%===============================================================================

%/////////
\begin{frame}[c,allowframebreaks,fragile]{Citations}
%
\begin{block}{\texttt{\textbackslash{}ccite} command---e.g., \ccite{bibtex-patashnik88}}
\begin{verbatim}
\ccite{bibtex-patashnik88}
\end{verbatim}
\begin{itemize}
	\item to be used instead of standard \texttt{\textbackslash{}cite} command
	\item light, for the dark surfaces a citation marker usually sits on
	\begin{itemize}
		\item as here above in the \texttt{block} header
	\end{itemize}
	\item can be used as a note\ccite{bibtex-patashnik88}, with no space before---faint in running text, by design
	\item since style 1.6.1 it takes the optional note of \texttt{\textbackslash{}cite}\ccite[p.~42]{bibtex-patashnik88}
	\begin{itemize}
		\item as here above, from \verb|\ccite[p.~42]{bibtex-patashnik88}|
	\end{itemize}
\end{itemize}
\end{block}
%
\begin{alertblock}{\texttt{\textbackslash{}cccite} command---e.g., \cccite{bibtex-patashnik88}}
\begin{verbatim}
\cccite{bibtex-patashnik88}
\end{verbatim}
\begin{itemize}
	\item the lighter of the two, for the loudest surfaces
	\begin{itemize}
		\item as here above in the \texttt{alertblock} header
	\end{itemize}
	\item can be used as a note with no space before, in the same way as \texttt{\textbackslash{}ccite}
	\item takes the optional note in the same way, as in \verb|\cccite[p.~42]{bibtex-patashnik88}|
	\begin{itemize}
		\item as below in the \texttt{alertblock} header
	\end{itemize}
\end{itemize}
\end{alertblock}
%
\begin{alertblock}{\texttt{\textbackslash{}cccite} with a note---e.g., \cccite[p.~42]{bibtex-patashnik88}}
\begin{verbatim}
\cccite[p.~42]{bibtex-patashnik88}
\end{verbatim}
\end{alertblock}
%
\end{frame}
%/////////

%/////////
\begin{frame}[c,fragile]{URLs}
%
\begin{block}{\texttt{\textbackslash{}uurl} command}
\begin{itemize}
	\item to be used instead of standard \texttt{\textbackslash{}url} command
	\item[e.g.] \verb|\uurl{https://apice.unibo.it}| prints as \uurl{https://apice.unibo.it}
	\item since style 1.6.2 no character in the address needs escaping
\end{itemize}
\end{block}
%
\begin{block}{\texttt{\textbackslash{}uuurl} command\textemdash{}e.g., \uuurl{https://apice.unibo.it}}
\begin{itemize}
	\item to be used instead of standard \texttt{\textbackslash{}url} command over dark backgrounds
	\item[e.g.] see \verb|\uuurl{https://apice.unibo.it}| above in this \texttt{block} header
\end{itemize}
\end{block}
%
\end{frame}
%/////////

%/////////
\begin{frame}[c,fragile]{DOI}
%
\begin{block}{\texttt{\textbackslash{}ddoi} command}
\begin{itemize}
	\item to be used for the styled DOI URLs
	\item[e.g.] \verb|\ddoi{10.1007/978-3-032-22940-3}| prints as \ddoi{10.1007/978-3-032-22940-3} % chktex 8
	\item[e.g.] since style 1.6.2 an underscore needs no escaping: \verb|\ddoi{10.1007/978-3-032-22940-3_12}| prints as \ddoi{10.1007/978-3-032-22940-3_12} % chktex 8
\end{itemize}
\end{block}
%
\begin{block}{\texttt{\textbackslash{}dddoi} command}
\begin{itemize}
	\item to be used for the styled DOI URLs
	\item[e.g.] \verb|\dddoi{10.1007/978-3-032-22940-3}| prints as \dddoi{10.1007/978-3-032-22940-3} % chktex 8
\end{itemize}
\end{block}
%
\end{frame}
%/////////

%/////////
\begin{frame}[c,fragile]{Alert}
%
\begin{block}{\texttt{\textbackslash{}aalert} command\textemdash{}e.g., \aalert{alerted text}}
\begin{itemize}
	\item to be used instead of standard \texttt{\textbackslash{}alert} command over dark backgrounds
	\item[e.g.] see \verb|\aalert{alerted text}| above in this \texttt{block} header
\end{itemize}
\end{block}
%
\end{frame}
%/////////

%/////////
\begin{frame}[c,fragile,allowframebreaks]{Speakers \emph{vs.} Authors}
%
\begin{exampleblock}{\texttt{\textbackslash{}speaker} command\textemdash{}e.g., \speaker{Diego Zorro}}
\begin{itemize}
	\item to be used within \texttt{\textbackslash{}author} standard Beamer command to single out the actual speaker among the authors
	\item[e.g.] as in
\begin{verbatim}
	\author[Garcia \and Zorro]
	{Sarg Garcia \and \speaker{Diego Zorro}}
\end{verbatim}
	\item and in the author specification of this template
\end{itemize}
\end{exampleblock}
%
\begin{block}{\texttt{\textbackslash{}sspeaker} command\textemdash{}e.g., \sspeaker{Diego Zorro}}
\begin{itemize}
	\item to be used within \texttt{\textbackslash{}author} standard Beamer command to single out the actual speaker among the authors in the short form
	\item[e.g.] as in
\begin{verbatim}
	\author[Garcia \and \sspeaker{Zorro}]
	{Sarg Garcia \and \speaker{Diego Zorro}}
\end{verbatim}
	\item and in the author specification of this template
\end{itemize}
\end{block}

\end{frame}
%/////////

%===============================================================================
\section{New Options}
%===============================================================================

%/////////
\begin{frame}[c,fragile,allowframebreaks]{\texttt{apice} option}
%
\begin{alertblock}{\texttt{apice} Beamer option}
\begin{itemize}
	\item the option \texttt{apice} can be used in the presentation preamble
	\begin{itemize}
		\item as in
		\begin{verbatim}
\documentclass[presentation,apice]{beamer}...
		\end{verbatim}
	\end{itemize}
	\item enables the display of the \BibTeX{} field \texttt{apice}
	\begin{itemize}
		\item containing the relative URL of the publication on the \alert{\href{https://apice.unibo.it/}{\textsf{APICe}}} XWiki
	\end{itemize}
	\item properly linked to the corresponding \href{https://apice.unibo.it/}{\textsf{APICe}} page
\end{itemize}
\end{alertblock}
%
\begin{exampleblock}{Example}
%
The following \BibTeX{} entry in the \texttt{.bib} file with the \texttt{apice} option enabled
%
{\tiny\begin{verbatim}
@manual{bibtex-patashnik88,
    apice = {BibtexPatashnik88},
    author = {Patashnik, Oren},
    booktitle = {\BibTeXing},
    organization = {CTAN, The Comprehensive TeX Archive Network},
    title = {\BibTeX{}ing},
    url = {https://mirrors.ctan.org/biblio/bibtex/base/btxdoc.pdf},
    year = 1988}
\end{verbatim}} %chktex 31
%
would translate into the following entry in the presentation bibliography
%
\begin{thebibliography}{}
\bibitem[Patashnik, 1988]{bibtex-patashnik1988}
Patashnik, O. (1988)
\newblock {\bfseries {\em \BibTeX{}ing}}
\newblock CTAN, The Comprehensive TeX Archive Network
\newblock\apicepubpar{BibtexPatashnik88} \uuurl{https://mirrors.ctan.org/biblio/bibtex/base/btxdoc.pdf}
\end{thebibliography}
%
\end{exampleblock}
%
\end{frame}
%/////////

%/////////
\begin{frame}[c,fragile]{\texttt{bologna}, \texttt{cesena} options}
%
\begin{alertblock}{Which seal watermarks the slides\textemdash{}none, unless asked for}
\begin{itemize}
	\item \texttt{bologna}\textemdash{}the seal of the \emph{Alma Mater Studiorum} as a whole
	\item \texttt{cesena}\textemdash{}the seal of the \emph{Campus di Cesena}
	\begin{itemize}
		\item the two are told apart by the border, solid for the Campus and an outline for the Ateneo
	\end{itemize}
	\item both sit in the lower-right corner, cropped by the page edge
	\begin{itemize}
		\item the one on this very slide is the one this deck selected
	\end{itemize}
	\item \alert{neither option, no seal}\textemdash{}the slides stay bare
	\begin{itemize}
		\item the seals are institutional marks, so a deck claims one out loud
		\item which leaves the style free to use for whoever does not belong\textemdash{}or no longer, or not yet belongs\textemdash{}to the \emph{Alma Mater}
	\end{itemize}
\end{itemize}
\end{alertblock}
%
\begin{exampleblock}{Example}
{\scriptsize\begin{verbatim}
\documentclass[presentation,bologna]{beamer}
    \mode<presentation>{\usetheme{AMSSilverBlue}}
\end{verbatim}} %chktex 31
\end{exampleblock}
%
\end{frame}
%/////////

%===============================================================================
\section*{} 
%===============================================================================

%/////////
\frame{\titlepage}
%/////////

%===============================================================================
\section*{\refname}
%===============================================================================

%%%%
\setbeamertemplate{page number in head/foot}{}
%/////////
\begin{frame}[c,noframenumbering]{\refname}
%\begin{frame}[t,allowframebreaks,noframenumbering]{\refname}
%	\tiny
	\scriptsize
%	\footnotesize
	\bibliographystyle{apalike-AMS}
	\bibliography{AMSSilverBlue-template}
\end{frame}
%/////////

%%%%%%%%%%%%%%%%%%%%%%%%%%%%%%%%%%%%%%%%%%%%%%%%%%%%%%%%%%%%%%%%%%%%%%%%%%%%%%%%
\end{document}
%%%%%%%%%%%%%%%%%%%%%%%%%%%%%%%%%%%%%%%%%%%%%%%%%%%%%%%%%%%%%%%%%%%%%%%%%%%%%%%%
